\SourcePage{378}
\SGAsection{5}{$\ell=p$ 情形的正交性定理（半稳定情形）；阿贝尔概形 $B_o$ 与 $B'_o$ 的对偶性；固定部分的刻画；好约化、本质好约化与半稳定约化判据}
\textbf{5.1。}置身于 2.2.3 的情形：$S$ 是 \Term{henselian}{Hensel} 的，$A^o$ 是 $\ell$ \Term{divisible}{可除}的，因而\Term{fixed-part}{固定部分}
\[
T_\ell(A^o)^f\lhook\joinrel\longrightarrow T_\ell(A)
\]
定义为由 $({}_{\ell^\nu}A^o)^f$ 组成的\Term{barsotti-tate-group}{Barsotti--Tate pro-$\ell$ 群}。现在我们要给出 $T_\ell(A^o)$ 的\Term{toric-part}{环面部分}（2.3.2）的一个定义，不再要求 2.3 中的假设 $\ell\ne p$，但改为假设 $S$ \Term{complete}{完备}\footnote{其实这个假设并非必要，见下文 6.1。}。为此，对每个整数 $n\ge0$，引入\Term{local-artinian-scheme}{局部 Artin 概形}
\[
S_n=\operatorname{Spec}(V/\mathfrak m^{n+1}),\tag{5.1.1}
\]
即 $s$ 在 $S$ 中的第 $n$ 个\Term{infinitesimal-neighborhood}{无穷小邻域}；并对任意 $S$-\Term{scheme}{概形} $X$ 记 $X_n=X\times_S S_n$。由于 $(A^o)_n=(A_n)^o$ 在 $S_n$ 上\Term{smooth}{光滑}，已知（SGA 3 IX 3.6 bis）$T_o$ 是唯一一个\Term{subtorus}{子环面}
\[
T_n\subset A_n^o\tag{5.1.2}
\]
的\Term{special-fiber}{特殊纤维}。由 $T_n$ 的唯一性，它们彼此通过\Term{base-change}{基变换} $S_{n'}\to S_n$ 得到。换言之，它们在 $S$ 上的\Term{formal-scheme}{形式概形}\Term{category}{范畴}中，定义了群 $\widehat A^o$ 的一个\Term{formal-subtorus}{形式子环面} $\widehat T$；这里 $\widehat A^o$ 是 $A^o$ 沿\Term{special-fiber}{特殊纤维}的\Term{formal-completion}{形式完备化}：
\[
\widehat T\subset\widehat A^o.\tag{5.1.3}
\]

\SourcePage{379}
由于\Term{functor}{函子} $X\mapsto\widehat X$（沿\Term{special-fiber}{特殊纤维}的\Term{formal-completion}{形式完备化}）给出 $S$ 上\Term{finite-scheme}{有限概形}的\Term{category}{范畴}与 $S$ 上\Term{finite}{有限}\Term{formal-scheme}{形式概形}的\Term{category}{范畴}之间的\Term{equivalence-of-categories}{等价}（EGA III 4.8），我们就把 $S$ 上的\Term{finite-scheme}{有限概形}与对应的\Term{formal-scheme}{形式概形}直接等同。在这一约定下，\Term{fixed-part}{固定部分}的定义立即给出，对每个 $\nu\ge0$，有
\[
({}_{\ell^\nu}A^o)^f={}_{\ell^\nu}(\widehat A^o),\tag{5.1.4}
\]
也可以把它写成\Term{projective-system}{投射系}的恒等式
\[
T_\ell(A^o)^f=T_\ell(\widehat A^o).\tag{5.1.5}
\]
因此，取 $T_\ell(\widehat T)\subset T_\ell(\widehat A^o)$，就得到 $T_\ell(A^o)$ 的\Term{fixed-part}{固定部分}的一个 \Term{barsotti-tate-group}{Barsotti--Tate 子 pro-群}。我们将 $T_\ell(A^o)$ 的\Term{toric-part}{环面部分}定义为
\[
T_\ell(A^o)^t\overset{\mathrm{def}}=T_\ell(\widehat T)
\subset T_\ell(\widehat A^o)=T_\ell(A^o)^f.\tag{5.1.6}
\]
于是，按定义有
\[
T_\ell(A^o)^t=({}_{\ell^\nu}\widehat T)_{\nu\ge0},\tag{5.1.7}
\]
其中\Term{finite}{有限}\Term{flat}{平坦} $S$-群 ${}_{\ell^\nu}\widehat T$ 本身可描述为\Term{finite}{有限}\Term{flat}{平坦} $S_n$-群
\[
{}_{\ell^\nu}(T_n),\qquad n\text{ 变化}\tag{5.1.8}
\]
的\Term{inductive-limit}{归纳极限}；也就是说，它的\Term{affine-algebra}{仿射代数}是群 (5.1.8) 的\Term{affine-algebra}{仿射代数}的\Term{projective-limit}{投射极限}。当然还有
\[
(T_\ell(A^o)^t)_n=T_\ell(T_n).\tag{5.1.9}
\]

\SourcePage{380}
当 $n=0$ 时，这正是公式 (2.3.3)：在 $\ell\ne p$ 的情形，我们曾用它刻画\Term{toric-part}{环面部分}（作为 pro-\Term{finite}{有限}\Term{etale}{平展}群）。这说明当 $\ell\ne p$ 时，定义 (5.1.6) 与 2.3 中的定义相容。

取\Term{generic-fiber}{一般纤维}，$T_\ell(A^o)$ 的\Term{fixed-part}{固定部分}和\Term{toric-part}{环面部分}给出 $T_\ell(A_K)$ 中的\Term{fixed-part}{固定部分}与\Term{toric-part}{环面部分}：
\[
T_\ell(A_K)^t\subset T_\ell(A_K)^f\subset T_\ell(A_K).\tag{5.1.10}
\]
这些都是 \Term{barsotti-tate-group}{Barsotti--Tate pro-$\ell$ 群}之间的包含。考虑与\Term{dual-abelian-scheme}{对偶阿贝尔概形} $A'_K$ 对应的类似\Term{filtration}{滤过}：
\[
T_\ell(A'_K)^t\subset T_\ell(A'_K)^f\subset T_\ell(A'_K).\tag{5.1.11}
\]
于是有下述结果，它推广了 2.4（至少在 $S$ \Term{complete}{完备}的情形如此，而 2.4 中很容易归结到这种情形）：
\begin{SGAresult}{\Term{theorem}{定理}}{5.2}
设 $S$ 为\Term{complete}{完备}\Term{trait}{离散赋值环谱}，$A_K$ 为 $K$ 上的\Term{abelian-scheme}{阿贝尔概形}，$A'_K$ 为\Term{dual-abelian-scheme}{对偶阿贝尔概形}；$\ell$ 为一个\Term{prime-number}{素数}，使\Term{connected}{连通} \Term{neron-model}{Néron 模型} $A^o$ 是 $\ell$ \Term{divisible}{可除}的（参见 2.2.1）。这样便有 pro-群 $T_\ell(A_K)$ 和 $T_\ell(A'_K)$ 的\Term{filtration}{滤过} (5.1.10) 与 (5.1.11)。在这些条件下，\Term{toric-part}{环面部分}可由\Term{fixed-part}{固定部分}通过公式
\[
T_\ell(A_K)^t=T_\ell(A_K)^f\cap(T_\ell(A'_K)^f)^\perp\tag{5.2.1}
\]
及交换 $A_K$ 与 $A'_K$ 的角色后所得的类似公式恢复；这里 $\perp$ 表示关于\Term{canonical-pairing}{典范配对} (1.0.2) 的\Term{orthogonal-complement}{正交补}。
\SourcePage{381}
此外，当 $A_K$ 在 $S$ 上具有\Term{semistable-reduction}{半稳定约化}时（$\ell=p$ 时其余假设已蕴含这一点），甚至有公式
\[
T_\ell(A_K)^t=(T_\ell(A'_K)^f)^\perp.\tag{5.2.2}
\]
\end{SGAresult}

\textbf{5.2.3。}若 $\ell$ 与 $K$ 的\Term{characteristic}{特征}互素，公式 (5.2.1) 的含义很清楚：像 2.4 那样，取有关 pro-群的 $\overline K$-点，便可用所得带 $\pi=\operatorname{Gal}(\overline K/K)$ \Term{action}{作用}的普通 $\mathbb Z_\ell$-模解释它。现在利用 \Term{barsotti-tate-group}{Barsotti--Tate pro-$\ell$ 群} $G=(G(\nu))$ 的\Term{dual}{对偶}
\[
D(G)=\operatorname{Hom}(G,T_\ell(\mathbb G_m))\tag{5.2.3.1}
\]
这一概念，说明\Term{generic}{一般}情形。按定义，它是 Cartier \Term{dual}{对偶} $D(G(\nu))$ 的\Term{projective-system}{投射系}（SGA 3 VII 3.3.1 及 VIII 1），其中过渡\Term{morphism}{态射}以显然的方式定义，例如由各 $G(\nu)$ 间的包含\Term{morphism}{态射}\Term{transpose}{转置}得到。\Term{abelian-variety}{阿贝尔簇}的\Term{duality}{对偶理论}告诉我们，5.2 所说的\Term{pairing}{配对}实际上定义了\Term{isomorphism}{同构}
\[
T_\ell(A'_K)\simeq D(T_\ell(A_K)),\tag{5.2.3.2}
\]
由复合便得到\Term{morphism}{态射}
\[
T_\ell(A_K)^f\longrightarrow D(T_\ell(A'_K)^f).\tag{5.2.3.3}
\]
公式 (5.2.1) 断言，上述\Term{homomorphism}{同态}的\Term{kernel}{核}（先验地是第一项的一个子 pro-群）恰好是\Term{toric-part}{环面部分} $T_\ell(A_K)^t$。应注意，这里说的是 \Term{pro-object}{pro-对象}的\Term{isomorphism}{同构}，而非\Term{projective-system}{投射系}的\Term{isomorphism}{同构}：我们并未断言，对每个整数 $\nu$，$T_\ell(A_K)^t(\nu)$ 等于
\SourcePage{382}
$T_\ell(A_K)^f(\nu)\to D(T_\ell(A'_K)^f)(\nu)$ 的\Term{kernel}{核}（这样的断言\Term{generic}{一般}是错误的）；后者先验地可能更大。不过，下文将看到，当 $A_K$ 具有\Term{semistable-reduction}{半稳定约化}时，后者并不更大；特别地，根据 5.2 的假设，$\ell=p$ 时即属此情形。

\textbf{5.2.4。}公式 (5.2.1) 的证明实际上就是几乎原封不动地重复 2.4 的证明。注意，产生 (5.2.3.3) 的诱导\Term{pairing}{配对}
\[
T_\ell(A_K)^f\times T_\ell(A'_K)^f\longrightarrow T_\ell(\mathbb G_{mK})\tag{5.2.4.1}
\]
来自 $S$ 上 \Term{barsotti-tate-group}{Barsotti--Tate pro-群}的\Term{pairing}{配对}
\[
T_\ell(A^o)^f\times T_\ell(A'^o)^f\longrightarrow T_\ell(\mathbb G_{mS}),\tag{5.2.4.2}
\]
而后者又由 pro-群的\Term{pairing}{配对}
\[
T_\ell(A^o)\times T_\ell(A'^o)\longrightarrow T_\ell(\mathbb G_{mS})\tag{5.2.4.3}
\]
诱导；它从 $(A^o,A'^o)$ 被 $\mathbb G_{mS}$ 的\Term{biextension}{双扩张} $W$ 得到。因此，由 (5.2.4.2) 取\Term{special-fiber}{特殊纤维}得到的\Term{pairing}{配对}
\[
T_\ell(A_o^o)\times T_\ell(A_o'^o)\longrightarrow T_\ell(\mathbb G_{mk})\tag{5.2.4.4}
\]
恰好就是由 $(A_o^o,A_o'^o)$ 被 $\mathbb G_{mk}$ 的\Term{biextension}{双扩张} $W_o$ 得到的\Term{pairing}{配对}。根据第 VIII 讲 4.11 和 1.5，\Term{pairing}{配对} (5.2.4.4) 的左\Term{annihilator}{零化子}等于 $T_\ell(T_o)$；这恰好是\Term{toric-part}{环面部分} $T_\ell(A^o)^t$ 的\Term{special-fiber}{特殊纤维}。因此，回到 5.2.3 所阐明的左\Term{annihilator}{零化子}定义，利用下述\Term{lemma}{引理}即可得出结论。

\SourcePage{383}
\begin{SGAresult}{\Term{lemma}{引理}}{5.2.5}
设 $u:G\to H$ 为 $S$ 上 \Term{barsotti-tate-group}{Barsotti--Tate pro-$\ell$ 群}的\Term{homomorphism}{同态}。若 $u_o:G_o\to H_o$ 为零，则 $u$ 为零。若 $u_o:G_o\to H_o$ 为\Term{isogeny}{同源}（相应地，为\Term{isomorphism}{同构}），则 $u$ 为\Term{isogeny}{同源}（相应地，为\Term{isomorphism}{同构}）。
\end{SGAresult}
事实上，对复合\Term{morphism}{态射}
\[
T_\ell(A^o)^t\longrightarrow T_\ell(A^o)^f\longrightarrow D(T_\ell(A'^o)^f)
\]
应用 5.2.5 的第一个断言，其中第二个\Term{morphism}{态射}来自\Term{pairing}{配对} (5.2.4.2)。可知此复合为零，因而所考虑的\Term{pairing}{配对}在左边的\Term{toric-part}{环面部分}上为零。由\Term{symmetry}{对称性}，它在右边的\Term{toric-part}{环面部分}上也为零，故来自商，也就是“\Term{abelian-part}{阿贝尔部分}”上的\Term{pairing}{配对}
\[
T_\ell(A^o)^{\mathrm{ab}}\times T_\ell(A'^o)^{\mathrm{ab}}\longrightarrow T_\ell(\mathbb G_m).
\]
容易核实，这些商仍是 \Term{barsotti-tate-group}{Barsotti--Tate pro-群}。剩下只需证明此\Term{pairing}{配对}对应的\Term{morphism}{态射}
\[
T_\ell(A^o)^{\mathrm{ab}}\longrightarrow D(T_\ell(A'^o)^{\mathrm{ab}})
\]
是\Term{isogeny}{同源}，这仍由 5.2.5 得出。

5.2.5 的证明很容易，留给读者；此结果在任意\Term{connected}{连通}\Term{noetherian}{Noether} 底\Term{scheme}{概形}上都成立。另外，在下文 7.4.7 中，通过稍微细致地分析 $(A_n^o,A_n'^o)$ 的\Term{biextension}{双扩张} $W_n$，我们会对刚才两次用到 5.2.5 的特殊情形，给出一个相当明显的直接证明。

\SourcePage{384}
\textbf{5.2.6。}现在假设 $A_K$ 在 $S$ 上具有\Term{semistable-reduction}{半稳定约化}；于是 $\ell$ 可取任意\Term{prime-number}{素数}（可以等于 $p$）。此外，由\Term{semistable-reduction}{半稳定约化}的 Galois 判据 3.5，当 $\ell\ne p$ 时有
\[
(T_\ell(A'_K)^f)^\perp\;(=T_\ell(A_K)^t)\subset T_\ell(A_K)^f.\tag{5.2.7}
\]
要完成 5.2 的证明，只需证明这个包含对 $\ell$ 无须任何\Term{restriction}{限制}。

事实上，左边\Term{canonical}{典范}\Term{isomorphism}{同构}于 $D(T_\ell(A'_K)/T_\ell(A'_K)^f)$，因此其“\Term{rank}{秩}”（或“\Term{height}{高度}”）等于 $2n-2\alpha-\mu=\mu$（采用 (2.1.9) 的记号，依假设 $\lambda=0$），这也正是 $T_\ell(A_K)^t$ 的\Term{rank}{秩}。由于包含
\[
T_\ell(A_K)^t\subset T_\ell(A_K)\tag{*}
\]
是\Term{strict-subgroup}{严格}的，意即它已在所考察\Term{projective-system}{投射系}的每一项上诱导\Term{monomorphism}{单态射}，故容易得到：在 $T_\ell(A_K)$ 的所有与它同\Term{rank}{秩}的 \Term{barsotti-tate-group}{Barsotti--Tate 子 pro-群}中，(*) 的第一项是极大的。因此 (5.2.7) 中的等式确实成立；由\Term{orthogonality-theorem}{正交性定理}，这\Term{equivalence}{等价}于 (5.2.7) 两个极端项之间的包含。

\textbf{5.3。}利用互为\Term{dual}{对偶}的 pro-群 $T_\ell(A_K)$ 与 $T_\ell(A'_K)$ 的\Term{strict-subgroup}{严格}（含义如上）\Term{barsotti-tate-group}{Barsotti--Tate 子 pro-群}之间的\Term{orthogonality}{正交}对应，由 (5.2.2) \Term{form}{形式}地得到：这些 pro-群之间的\Term{pairing}{配对} (1.0.2) 诱导一个\Term{pairing}{配对}
\[
T_\ell(A_K)^{\mathrm{ab}}\times T_\ell(A'_K)^{\mathrm{ab}}\longrightarrow T_\ell(\mathbb G_{mK}),\tag{5.3.2}
\]
\SourcePage{385}
它是\Term{perfect-duality}{完美对偶}，即把左边的两个因子分别与另一因子的\Term{dual}{对偶}等同（见 5.2.2）。注意，上标 $\mathrm{ab}$ 表示“\Term{abelian-part}{阿贝尔部分}”，也就是\Term{fixed-part}{固定部分}模去\Term{toric-part}{环面部分}的商。另一方面，(5.3.2) 显然由\Term{pairing}{配对}
\[
T_\ell(A^o)^{\mathrm{ab}}\times T_\ell(A'^o)^{\mathrm{ab}}\longrightarrow T_\ell(\mathbb G_{mS})\tag{5.3.3}
\]
诱导；后者由\Term{pairing}{配对} (5.2.4.2) 取商得到。它在\Term{generic-fiber}{一般纤维}上是\Term{perfect-duality}{完美对偶}，即对应的\Term{homomorphism}{同态}
\[
T_\ell(A^o)^{\mathrm{ab}}\longrightarrow D(T_\ell(A'^o)^{\mathrm{ab}})\tag{5.3.4}
\]
在\Term{generic-fiber}{一般纤维}上为\Term{isomorphism}{同构}，于是自然希望\Term{pushout}{推出}它在整个 $S$ 上，\Term{equivalence}{等价}地，在\Term{special-fiber}{特殊纤维}上也是如此。当 $\ell\ne p$ 时，这个推论显然成立；当 $\ell=p$ 时，至少在\Term{mixed-characteristic}{混合特征}情形，由 Tate 的一个深刻\Term{theorem}{定理} [33，\Term{theorem}{定理} 4]，它也成立。不过，下文 7.4.6 将直接建立这一点，无须借助 Tate \Term{theorem}{定理}。

为说明我们要建立的结果的几何意义，注意，由 (5.3.3) 取\Term{special-fiber}{特殊纤维}得到的\Term{pairing}{配对}，恰好是由 $(A_o^o,A_o'^o)$ 的\Term{biextension}{双扩张} $W_o$ 取商得到的\Term{pairing}{配对}
\[
T_\ell(B_o)\times T_\ell(B'_o)\longrightarrow T_\ell(\mathbb G_{mk}),\tag{5.3.5}
\]
\SourcePage{386}
其中 $B_o$ 和 $B'_o$ 分别表示 $A_o^o$ 和 $A_o'^o$ 的\Term{abelian-part}{阿贝尔部分}。由第 VIII 讲 4.8，\Term{biextension}{双扩张} $W_o$ 来自 $(B_o,B'_o)$ 被 $\mathbb G_{mk}$ 的\Term{biextension}{双扩张} $W_o^{\mathrm{ab}}$，后者在唯一\Term{isomorphism}{同构}的意义下确定。因此，\Term{pairing}{配对} (5.3.5) 也\Term{solvable}{可解}释为从\Term{divisorial-correspondence}{除子对应} $W_o^{\mathrm{ab}}$ 得到的\Term{pairing}{配对}；对应的\Term{homomorphism}{同态}
\[
T_\ell(B_o)\longrightarrow D(T_\ell(B'_o))\simeq T_\ell(D(B'_o))\tag{5.3.6}
\]
（最后一项中的字母 $D$ 表示\Term{dual-abelian-scheme}{对偶阿贝尔概形}）恰好由 $W_o^{\mathrm{ab}}$ 定义的\Term{homomorphism}{同态}
\[
B_o\longrightarrow D(B'_o)\tag{5.3.7}
\]
经 $T_\ell$ 的\Term{functoriality}{函子性}得到。于是，由熟知的对应关系（参见 [5]），(5.3.5) 是\Term{perfect-duality}{完美对偶}，也意味着 (5.3.7) 的\Term{kernel}{核}的 $\ell$-\Term{primary-component}{准素部分}为零；根据 1.5，这个\Term{kernel}{核}是 $k$ 上的\Term{finite-group}{有限群}。对所有 $\ell$ 都如此，就意味着\Term{isogeny}{同源} (5.3.7) 是\Term{isomorphism}{同构}，即\Term{divisorial-correspondence}{除子对应} $W_o^{\mathrm{ab}}$ 是 $B_o$ 与 $B'_o$ 之间的\Term{perfect-duality}{完美对偶}。前面的讨论仅能说明 (5.3.7) 的\Term{kernel}{核}是一个 $p$-群。现在陈述更完整的结果：
\begin{SGAresult}{\Term{theorem}{定理}}{5.4}
设 $S$ 为\Term{trait}{离散赋值环谱}，$A_K$ 为 $S$ 上的\Term{abelian-scheme}{阿贝尔概形}，且在 $S$ 上具有\Term{semistable-reduction}{半稳定约化}。分别以 $B_o,B'_o$ 表示 $A_K$ 及\Term{dual-abelian-scheme}{对偶阿贝尔概形} $A'_K$ 的 \Term{neron-model}{Néron 模型}的\Term{abelian-part}{阿贝尔部分}。那么，$(B_o,B'_o)$ 上的\Term{canonical-divisorial-correspondence}{典范除子对应} $W_o^{\mathrm{ab}}$，即由第 VIII 讲 4.8 从 $(A_o^o,A_o'^o)$ 的\Term{canonical-biextension}{典范双扩张} $W_o$（1.4 的\Term{biextension}{双扩张} $W$ 的\Term{special-fiber}{特殊纤维}）得到的对应，是 $B_o$ 与 $B'_o$ 之间的\Term{perfect-duality}{完美对偶}；
\SourcePage{387}
也就是说，它定义了这两个\Term{abelian-scheme}{阿贝尔概形}中的每一个与另一个的\Term{dual}{对偶}之间的\Term{isomorphism}{同构}。
\end{SGAresult}
5.4 在\Term{generic}{一般}情形下的证明将在第 7 节给出，作为 Raynaud 的一个十分自然且重要的构造的推论。这里只须指出，由 3.3，证明 5.4 时可以假设 $S$ \Term{complete}{完备}。

\textbf{5.5。}为便于表达\Term{semistable-reduction}{半稳定约化}情形中\Term{filtration}{滤过} (5.1.10) 和 (5.1.11) 的两个极端因子，引入\Term{twisted-constant-group-scheme}{扭常值群概形}
\[
\underline M'_o=D(T_o),\qquad \underline M_o=D(T'_o),\tag{5.5.1}
\]
它们是 $A,A'$ 的\Term{special-fiber}{特殊纤维}的极大\Term{torus}{环面} $T_o,T'_o$ 的\Term{dual}{对偶}（SGA 3 X 5.7）。由于 $S$ 是\Term{complete}{完备}\Term{local-scheme}{局部概形}，可将这些\Term{group-scheme}{群概形}看成 $S$ 上\Term{twisted-constant-group-scheme}{扭常值群概形}
\[
\underline M',\ \underline M\quad\text{（在 $S$ 上）}\tag{5.5.2}
\]
的\Term{special-fiber}{特殊纤维}。例如，若\Term{torus}{环面} $T_o$ \Term{split}{分裂}（如 $k$ \Term{separably-closed}{可分闭}），即 $\underline M'_o$ 是普通群
\[
M'=\operatorname{Hom}_{k\text{-gr}}(T_o,\mathbb G_{mk})\tag{5.5.3}
\]
所定义的\Term{constant-group}{常值群}，那么 $\underline M'$ 就是 $M'$ 所定义的\Term{constant-group}{常值群}：
\[
\underline M'=M'_S.
\]
还可注意，由于 $T_o$ 与 $T'_o$ \Term{isogenous}{同源}（2.2.6），若 $T_o$ \Term{split}{分裂}，则 $T'_o$ 也\Term{split}{分裂}，因此在这种情形下还应引入
\SourcePage{388}
\[
M=\operatorname{Hom}_{k\text{-gr}}(T'_o,\mathbb G_{mk}),\tag{5.5.4}
\]
于是有 $\underline M=M_S$。这样，对每个整数 $n$，显然有\Term{isomorphism}{同构}
\[
\underline M'_n\simeq D(T_n),\qquad \underline M_n\simeq D(T'_n),\tag{5.5.5}
\]
这里由 $A'$ 定义 $T'_n$ 的方式，与由 $A$ 定义 $T_n$ 的方式相同（5.1.2）。由此得到\Term{canonical-isomorphism}{典范同构}（记 $\mathbb Z_\ell(1)=T_\ell(\mathbb G_{mS})$）：
\[
T_\ell(T_n)\simeq\underline M_n'{}^\vee\otimes_{\mathbb Z}\mathbb Z_\ell(1),\qquad
T_\ell(T'_n)\simeq\underline M_n^\vee\otimes_{\mathbb Z}\mathbb Z_\ell(1).\tag{5.5.6}
\]
再对 $n$ 取\Term{limit}{极限}，得到\Term{canonical-isomorphism}{典范同构}
\[
\begin{cases}
T_\ell(A^o)^t=T_\ell(\widehat T)\simeq\underline M'{}^\vee\otimes\mathbb Z_\ell(1),\\
T_\ell(A'^o)^t=T_\ell(\widehat T')\simeq\underline M^\vee\otimes\mathbb Z_\ell(1).
\end{cases}\tag{5.5.7}
\]
在\Term{semistable-reduction}{半稳定约化}情形中，利用\Term{orthogonality-theorem}{正交性定理} 5.2 以及关系 (5.2.2)，由\Term{duality}{对偶性}得到\Term{canonical-isomorphism}{典范同构}
\[
\begin{cases}
T_\ell(A_K)/T_\ell(A_K)^f\simeq\underline M_K\otimes_{\mathbb Z}\mathbb Z_\ell,\\
T_\ell(A')/T_\ell(A'_K)^f\simeq\underline M'_K\otimes_{\mathbb Z}\mathbb Z_\ell.
\end{cases}\tag{5.5.8}
\]
同样，由同一公式 (5.2.2)，得到\Term{canonical-isomorphism}{典范同构}
\[
\begin{cases}
T_\ell(A_K)/T_\ell(A_K)^t\simeq D(T_\ell(A'_K)^f),\\
T_\ell(A'_K)/T_\ell(A'_K)^t\simeq D(T_\ell(A_K)^f).
\end{cases}\tag{5.5.9}
\]

\SourcePage{389}
须记住，右端出现的群 $T_\ell(A_K)^f$ 和 $T_\ell(A'_K)^f$，分别是定义在 $S$ 上的 \Term{barsotti-tate-group}{Barsotti--Tate pro-$\ell$ 群} $T_\ell(A^o)^f$ 与 $T_\ell(A'^o)^f$ 在 $\eta$ 上的\Term{restriction}{限制}；这里 $D$ 表示 \Term{barsotti-tate-group}{Barsotti--Tate pro-群}意义下的\Term{dual}{对偶}，即 $\operatorname{Hom}(-,\mathbb Z_\ell(1))$。因此有：
\begin{SGAresult}{命题}{5.6}
假设 $S$ \Term{complete}{完备}，且 $A_K$ 在 $S$ 上具有\Term{semistable-reduction}{半稳定约化}。那么，对任意\Term{prime-number}{素数} $\ell$（可以等于 $S$ 的\Term{residue-characteristic}{剩余特征}），\Term{barsotti-tate-group}{Barsotti--Tate pro-$\ell$ 群} $T_\ell(A_K)^t$ 是 $S$ 上一个“\Term{toric-type}{环面型}”\Term{barsotti-tate-group}{Barsotti--Tate pro-群}的\Term{generic-fiber}{一般纤维}，后者由 (5.5.7) 给出；而 $T_\ell(A_K)/T_\ell(A_K)^f$ 是 $S$ 上一个 pro-\Term{etale}{平展}型 \Term{barsotti-tate-group}{Barsotti--Tate pro-群}的\Term{generic-fiber}{一般纤维}，后者由 (5.5.8) 给出。最后，群 $T_\ell(A_K)^f$ 与 $T_\ell(A_K)/T_\ell(A_K)^t$ 是定义在 $S$ 上的 \Term{barsotti-tate-group}{Barsotti--Tate pro-群}的\Term{generic-fiber}{一般纤维}，后两者分别等于 $T_\ell(A^o)^f$ 和
\[
D(T_\ell(A'^o)^f)=\operatorname{Hom}(T_\ell(A'^o)^f,\mathbb Z_\ell(1)).
\]
\end{SGAresult}
\textbf{5.6.1。}注意，对一个\Term{etale}{平展}\Term{covering}{覆盖} $S'$ 而言，使 $T_\ell(A_K)^t$ \Term{split}{分裂}与使 $T_\ell(A_K)/T_\ell(A_K)^f$ \Term{split}{分裂}\Term{equivalence}{等价}，因为这分别意味着 $S'$ 使\Term{twisted-constant-group-scheme}{扭常值群} $\underline M$ 与 $\underline M'$ \Term{split}{分裂}；而这两个群\Term{isogenous}{同源}，因为 $T_o$ 与 $T'_o$ \Term{isogenous}{同源}。同时还可看出，存在一个最小的\Term{etale}{平展}\Term{extension}{扩张} $S'$ 满足这一要求。

再指出下面这个由 2.2.3 的定义容易得到的推论。

\SourcePage{390}
\begin{SGAresult}{命题}{5.7}
设 $A^o$ 为 $\ell$ \Term{divisible}{可除}的，$S$ 为 \Term{henselian}{Hensel} 的。令 $F_K$ 为 $K$ 上一个“\Term{unramified}{非分歧}”的\Term{constructible-sheaf}{可构造} $\ell$ 进\Term{sheaf}{层}，即它可延拓成 $S$ 上的\Term{twisted-constant-sheaf}{扭常值} $\ell$ 进\Term{sheaf}{层} $F$；又令 $u:F_K\to T_\ell(A_K)$ 为 pro-群的\Term{homomorphism}{同态}。则 $u$ 可通过一个从 $F$ 到 $T_\ell(A_K)^f$ 的\Term{homomorphism}{同态}\Term{resolution}{分解}。
\end{SGAresult}
事实上，有
\[
u=(u(\nu))_{\nu\ge0},\qquad
u(\nu):F_K(\nu)\longrightarrow T_\ell(A_K)(\nu)={}_{\ell^\nu}A_K.\tag{5.7.1}
\]
由 $A$ 的\Term{universal-property}{泛性质}，由于 $F(\nu)$ 在 $S$ 上\Term{etale}{平展}从而\Term{smooth}{光滑}，\Term{morphism}{态射} $u(\nu)$ 来自一个\Term{homomorphism}{同态}
\[
v(\nu):F(\nu)\longrightarrow{}_{\ell^\nu}A.
\]
由\Term{fixed-part}{固定部分}的定义，此\Term{homomorphism}{同态}\Term{resolution}{分解}为
\[
v(\nu):F(\nu)\longrightarrow({}_{\ell^\nu}A)^f.
\]
由于它在\Term{special-fiber}{特殊纤维}上诱导的\Term{morphism}{态射}的像落在 $A_o$ 的无限 $\ell$ \Term{divisible}{可除}部分中，故 $v(\nu)$ 甚至通过 $({}_{\ell^\nu}A^o)^f$ \Term{resolution}{分解}。对 $\nu$ 取\Term{limit}{极限}，得到\Term{homomorphism}{同态}
\[
v:F\longrightarrow T_\ell(A^o)^f;
\]
它在\Term{generic-fiber}{一般纤维}上诱导的\Term{morphism}{态射}显然给出所求的\Term{resolution}{分解}。

\textbf{注记 5.7.2。}当 $\ell\ne p$ 时，$T_\ell(A_K)$ 的子 pro-群 $T_\ell(A_K)^f$ 的性质 5.7，显然唯一地刻画它为 $T_\ell(A_K)$ 中最大的“\Term{unramified}{非分歧}”$\ell$ 进\Term{subsheaf}{子层}。当 $\ell=p$ 时，命题 5.7 已不大合理；我们将给出一个远强于 5.7 所述的性质，
\SourcePage{391}
它能作为 $T_p(A_K)$ 的\Term{fixed-part}{固定部分}的内在刻画，仅用 $T_p(A_K)$ 的 \Term{barsotti-tate-group}{Barsotti--Tate pro-群}结构来表述——至少在\Term{mixed-characteristic}{混合特征}情形如此。

\begin{SGAresult}{\Term{theorem}{定理}}{5.8（M. Raynaud）}
设 $A_K$ 为 \Term{henselian}{Hensel} \Term{trait}{离散赋值环谱} $S$ 的\Term{function-field}{函数域} $K$ 上的\Term{abelian-scheme}{阿贝尔概形}，具有\Term{semistable-reduction}{半稳定约化}。令 $\ell$ 为一个\Term{prime-number}{素数}（有趣的情形是 $\ell=p$），考虑 \Term{barsotti-tate-group}{Barsotti--Tate pro-$\ell$ 群} $X_\eta=T_\ell(A_K)$ 及其“\Term{fixed-part}{固定部分}”$X_\eta^f$；后者是 \Term{barsotti-tate-group}{Barsotti--Tate pro-$\ell$ 群} $X^f=T_\ell(A^o)^f$ 的\Term{generic-fiber}{一般纤维}。当 $\ell=p$ 时，假设 $\operatorname{char}K=0$（因而在 $S$ 上可以使用 Tate \Term{theorem}{定理} [33，\Term{theorem}{定理} 4]）。设 $F$ 为 $S$ 上的 \Term{barsotti-tate-group}{Barsotti--Tate pro-$\ell$ 群}。则每个\Term{homomorphism}{同态} $F_\eta\to X_\eta$ 都通过 $X_\eta^f$ \Term{resolution}{分解}。
\end{SGAresult}
改取所考察的 \Term{barsotti-tate-group}{Barsotti--Tate pro-$\ell$ 群}对应的 ind-$\ell$ 群，由所引 Tate \Term{theorem}{定理}，上述命题\Term{equivalence}{等价}于：
\begin{SGAresult}{推论}{5.9}
若 $F$ 为 \Term{barsotti-tate-group}{Barsotti--Tate ind-$\ell$ 群}，则每个\Term{homomorphism}{同态} $u_\eta:F_\eta\to A_\eta$ 都唯一延拓为\Term{homomorphism}{同态} $u:F\to A$。
\end{SGAresult}
换言之，刻画 \Term{neron-model}{Néron 模型}的“Néron 性质”(1.1.2)，对 \Term{barsotti-tate-group}{Barsotti--Tate ind-群概形}的\Term{homomorphism}{同态}仍然成立，无须假设该 ind-\Term{group-scheme}{群概形}\Term{smooth}{光滑}。（只有 ind-\Term{etale}{平展}群的情形能直接由定义得到。）

\SourcePage{392}
\textbf{注记 5.9.1。}事实上，M. Raynaud 证明 5.9 时，并不要求 $A_K$ 具有\Term{semistable-reduction}{半稳定约化}；甚至对 $S$ 上任意 Néron \Term{group-scheme}{群概形} $A$ 都成立（其\Term{generic-fiber}{一般纤维}不必是\Term{abelian-scheme}{阿贝尔概形}）。当然，表述 5.9 时也无须要求 $S$ 为 \Term{henselian}{Hensel} 的。Raynaud 对这个更\Term{generic}{一般}情形的证明，是利用 \Term{neron-model}{Néron 模型}构造中的一个未发表的性质，归结到 5.8 所考察的特殊情形。这里我们只给出 Raynaud 对 5.9 的证明。

\begin{proof}[5.9 的证明]
已知 [33]，$F$ 是一个 ind-\Term{etale}{平展}的 \Term{barsotti-tate-group}{Barsotti--Tate 群} $F'$ 被一个\Term{special-fiber}{特殊纤维}\Term{connected}{连通}的 \Term{barsotti-tate-group}{Barsotti--Tate 群} $F''$ 的\Term{group-extension}{扩张}。因此，对每个整数 $n$，有\Term{exact-sequence}{正合列}
\[
0\longrightarrow F''(n)\longrightarrow F(n)\longrightarrow F'(n)\longrightarrow0,
\]
其中 $F'(n)$ \Term{connected}{连通}，$F''(n)_o$ 是\Term{purely-infinitesimal}{纯无穷小}的。给定的\Term{homomorphism}{同态} $u_\eta$ 由一列\Term{homomorphism}{同态}
\[
u_\eta(n):F(n)_\eta\longrightarrow A_\eta
\]
定义；只须分别把它们延拓为 $u(n):F(n)\to A$。下述\Term{lemma}{引理} 5.9.2 将问题归结为证明：由 $u_\eta$ 诱导的\Term{homomorphism}{同态} $u''_\eta:F''_\eta\to A_\eta$ 可延拓成 $u'':F''\to A$。因此可以假设 $F$ 的\Term{special-fiber}{特殊纤维}\Term{connected}{连通}。

我们可以把 $u_\eta$ 看成从 $F_\eta$ 到 $X_\eta=T_\ell(A_K)$ 所对应的 \Term{barsotti-tate-group}{Barsotti--Tate ind-$\ell$ 群} $Y_\eta$ 的\Term{homomorphism}{同态}，进而考虑复合
\SourcePage{393}
\[
v_\eta:F_\eta\longrightarrow Y_\eta\longrightarrow Y_\eta/Y_\eta^f.\tag{*}
\]
由 5.6，最后这个群可延拓为 $S$ 上的 ind-\Term{etale}{平展} \Term{barsotti-tate-group}{Barsotti--Tate 群} $Z$，即 $\underline M\otimes_{\mathbb Z}(\mathbb Q_\ell/\mathbb Z_\ell)$。由 5.8 中所引 Tate \Term{theorem}{定理}，\Term{homomorphism}{同态} (*) 可延拓为\Term{homomorphism}{同态}
\[
v:F\longrightarrow Z.
\]
由于 $F$ 的\Term{special-fiber}{特殊纤维}\Term{connected}{连通}，而 $Z$ 是 ind-\Term{etale}{平展}的，此\Term{morphism}{态射}为零。因此 $v_\eta$ 也为零，故 $u_\eta$ 通过 $Y_\eta^f$ \Term{resolution}{分解}。这就是所要证明的！
\end{proof}
现在只须证明：
\begin{SGAresult}{\Term{lemma}{引理}}{5.9.2}
设 $S$ 为\Term{trait}{离散赋值环谱}，$A$ 为 $S$ 上的 Néron \Term{group-scheme}{群概形}（1.1.2），
\[
1\longrightarrow G'\longrightarrow G\longrightarrow G''\longrightarrow1
\]
是 $S$ 上\Term{group-scheme}{群概形}的\Term{group-extension}{扩张}，其中 $G''$ \Term{smooth}{光滑}；又设
\[
u:G'\longrightarrow A
\]
为群\Term{homomorphism}{同态}，而
\[
v_\eta:G_\eta\longrightarrow A_\eta
\]
为延拓 $u_\eta$ 的群\Term{homomorphism}{同态}。假设由 $u$ \Term{pushout}{推出}的、$G''$ 被 $A$ 的\Term{group-extension}{扩张}\Term{representable}{可表}\footnote{M. Raynaud 未发表的结果实际上意味着，这个条件总是满足。}。例如，若 $G\to G''$ 对“\Term{finite-locally-free}{有限局部自由}”\Term{topology}{拓扑}\Term{local-scheme}{局部}地有\Term{section}{截面}（参见 SGA 3 IV 6.3），且 $A$ 在 $S$ 上\Term{quasi-projective}{拟射影}，则此条件满足。那么，存在唯一的群\Term{morphism}{态射}
\[
v:G\longrightarrow A
\]
同时延拓 $v_\eta$ 与 $u$。
\end{SGAresult}
\SourcePage{394}
\begin{proof}
事实上，最后的假设将我们归结到 $G'=A$、$u=\mathrm{id}$ 的情形；于是 $v_\eta:G_\eta\to A_\eta$ 是一个\Term{projection}{投影}，我们要将其延拓成\Term{projection}{投影} $v:G\to A$。但 $G$ 是两个\Term{smooth}{光滑}\Term{group-scheme}{群概形}的\Term{group-extension}{扩张}，所以\Term{smooth}{光滑}，因而结论由 $A$ 的 Néron 性质的定义得到。
\end{proof}

\begin{SGAresult}{推论}{5.10}
设 $A_K$ 为\Term{trait}{离散赋值环谱} $S$ 的\Term{function-field}{函数域} $K$ 上的\Term{abelian-scheme}{阿贝尔概形}，$\ell$ 为一个\Term{prime-number}{素数}。当 $\ell=p$ 时，假设 $\operatorname{char}K=0$，以便能够使用 5.8 中已经引用过的 Tate \Term{theorem}{定理}。$A_K$ 在 $S$ 上具有\Term{good-reduction}{好约化}的充要条件，是 \Term{barsotti-tate-group}{Barsotti--Tate pro-$\ell$ 群} $T_\ell(A_K)$ “在 $S$ 上具有\Term{good-reduction}{好约化}”，即它能延拓为 $S$ 上的 \Term{barsotti-tate-group}{Barsotti--Tate 群}。
\end{SGAresult}
（这正是 [8，4.9] 中猜想的结果。）
\begin{proof}
必要性显然，只须证明充分性；顺便指出，$\ell\ne p$ 的情形已经处理过（2.2.9）。可以假设 $S$ 为 \Term{henselian}{Hensel} 的（2.2.9）。先设 $A_K$ 在 $S$ 上具有\Term{semistable-reduction}{半稳定约化}。对 $F=X$ 应用 5.8，得到 $T_\ell(A_K)=T_\ell(A_K)^f$；这也意味着 $A_o^o$ 的极大\Term{torus}{环面}为零，即 $A_o^o$ 是\Term{abelian-scheme}{阿贝尔概形}，因而结论由 2.2.9 (iii$^o$) 得到。现在处理\Term{generic}{一般}情形。根据 2.2.9，显然可以假设 $S$ \Term{strictly-local}{严格局部}。取一个\Term{finite}{有限} Galois \Term{extension}{扩张} $K'$，使 $A_{K'}$ 在 $S$ 于 $K'$ 中的\Term{normalization}{正规化} $S'$ 上具有\Term{semistable-reduction}{半稳定约化}（3.6）。若以 $A^*$ 表示对应的 \Term{neron-model}{Néron 模型}，则 $G=\operatorname{Gal}(K'/K)$ \Term{action}{作用}在 $(A^*,S')$ 上，因而也\Term{action}{作用}在
\SourcePage{395}
$A_o^{*o}$ 上。根据 4.1 所见，只需证明 $G$ 在后者上\Term{trivial-action}{平凡作用}；这将蕴含 $A_K$ 具有\Term{semistable-reduction}{半稳定约化}，即回到已经处理的情形。令 $X$ 为 $S$ 上延拓 $X_\eta=T_\ell(A_K)$ 的 \Term{barsotti-tate-group}{Barsotti--Tate 群}。由已引用过的 Tate 唯一性\Term{theorem}{定理}，存在一个（唯一的）\Term{isomorphism}{同构}
\[
X\times_S S'\simeq T_\ell(A^{*o}),
\]
它在\Term{generic-fiber}{一般纤维}上诱导恒等。由\Term{transport-of-structure}{结构搬运}，此\Term{isomorphism}{同构}与 $G$ 的\Term{action}{作用}相容。而 $X$ 在 $X\times_S S'$ 的\Term{special-fiber}{特殊纤维}上的\Term{action}{作用}显然平凡，故它在 $T_\ell(A_o^{*o})$ 上\Term{trivial-action}{平凡作用}，从而显然在 $A_o^{*o}$ 上也\Term{trivial-action}{平凡作用}。证明完成。
\end{proof}

\textbf{注记 5.11。}一旦 Tate \Term{theorem}{定理}无须对 $\operatorname{char}K$ 作\Term{restriction}{限制}性假设就能证明，上述结果 5.8、5.9、5.10 及下文的 5.13 显然也就不再需要这种\Term{restriction}{限制}。

下面给出一个在别处也可能有用的定义。
\begin{SGAresult}{定义}{5.12}
设 $S$ 为 \Term{henselian}{Hensel} \Term{trait}{离散赋值环谱}，$\ell$ 为一个\Term{prime-number}{素数}，$X_\eta$ 为 $S$ 的\Term{generic-point}{一般点} $\eta$ 上的 \Term{barsotti-tate-group}{Barsotti--Tate pro-$\ell$ 群}。若 $X$ 能延拓为 $S$ 上的 \Term{barsotti-tate-group}{Barsotti--Tate 群}（参见 5.10），则称 $X$ 在 $S$ 上具有\Term{good-reduction}{好约化}。若存在 $X$ 的一个由\Term{strict-subgroup}{严格} \Term{barsotti-tate-group}{Barsotti--Tate 子群}\footnote{即作为\Term{projective-system}{投射系}也是\Term{subobject}{子对象}，也就是说逐项都是\Term{subobject}{子对象}。}组成的\Term{filtration}{滤过} $(\operatorname{Fil}^i(X_\eta))_{i\ge0}$，使 \Term{barsotti-tate-group}{Barsotti--Tate 群}
\[
\operatorname{Gr}^i(X_\eta)=\operatorname{Fil}^i(X_\eta)/\operatorname{Fil}^{i+1}(X_\eta)
\]
在 $S$ 上具有\Term{good-reduction}{好约化}，则称 $X$ 在 $S$ 上具有\Term{virtual-good-reduction}{虚好约化}（bonne réduction virtuelle）。若能选取\Term{length}{长度} $\le n$ 的\Term{filtration}{滤过}，即各 $\operatorname{Gr}^i$ 中至多有 $n$ 个非零，
\SourcePage{396}
则称 $X$ 在 $S$ 上具有\Term{virtual-good-reduction}{阶数为 $n$ 的虚好约化}。最后，若存在\Term{finite}{有限}\Term{extension}{扩张} $K'$，使 $X_{K'}$ 在 $S$ 于 $K'$ 中的\Term{normalization}{正规化} $S'$ 上具有\Term{good-reduction}{好约化}（相应地，\Term{virtual-good-reduction}{虚好约化}，或\Term{virtual-good-reduction}{阶数为 $n$ 的虚好约化}），则称 $X$ 在 $S$ 上具有\Term{potential-good-reduction}{本质好约化}（相应地，\Term{potential-virtual-good-reduction}{本质虚好约化}，或\Term{potential-virtual-good-reduction}{阶数为 $n$ 的本质虚好约化}）。
\end{SGAresult}

\textbf{5.12.1。}当 $\ell\ne p$ 时，$\eta$ 上（相应地，$S$ 上）的 \Term{barsotti-tate-group}{Barsotti--Tate pro-$\ell$ 群}\Term{category}{范畴}可用 $\pi=\pi_1(\eta)$（相应地，$\pi_o=\pi_1(S)$）的 $\ell$ 进\Term{representation}{表示}解释。于是立即看出，新引入的六个概念，\Term{separable}{可分}别表述为\Term{inertia-group}{惯性群} $I$ 在 $X_\eta(\overline K)$ 上的自然\Term{representation}{表示}具有以下性质：平凡、\Term{unipotent}{幺幂}、阶数为 $n$ 的\Term{unipotent}{幺幂}、本质平凡、本质\Term{unipotent}{幺幂}、阶数为 $n$ 的本质\Term{unipotent}{幺幂}。

\textbf{5.12.2。}设取 $X_\eta=T_\ell(A_K)$，其中 $A_K$ 为 $K$ 上的\Term{abelian-scheme}{阿贝尔概形}。5.6 已说明：若 $A_K$ 在 $S$ 上具有\Term{semistable-reduction}{半稳定约化}，则 $X_\eta$ 具有\Term{virtual-good-reduction}{阶数为 $2$ 的虚好约化}（稍后还会看到逆命题）；因此，由\Term{semistable-reduction}{半稳定约化}\Term{theorem}{定理} 3.6，$X_\eta$ 具有\Term{potential-virtual-good-reduction}{阶数为 $2$ 的本质虚好约化}。对 $K$ 上\Term{proper}{固有}\Term{smooth}{光滑}\Term{scheme}{概形}的 $H^1$ 而言，应把这个命题视为第 III 讲\SGAsourcegap{原印在此处的“III”之后留有空白，缺具体交叉引注}的几何\Term{monodromy}{单值性}\Term{theorem}{定理}的“正确”表述，它不排除 $\ell=p$ 的情形（原引文并未讨论这种情形）。对于高次 $H^i$，也应把
\SourcePage{397}
第 III 讲\SGAsourcegap{原印“donné dans III”之后留有空白，缺具体交叉引注}给出的几何\Term{monodromy}{单值性}\Term{theorem}{定理}（更确切地说，目前还是猜想）重新表述，以纳入“$p$ 进”系数，正如 0.4 所说。

下面概括 $K$ 上\Term{abelian-scheme}{阿贝尔概形}的\Term{reduced}{约化}性质与相应 \Term{barsotti-tate-group}{Barsotti--Tate 群}的\Term{reduced}{约化}性质之间的关系。
\begin{SGAresult}{命题}{5.13}
设 $S$ 为 \Term{henselian}{Hensel} \Term{trait}{离散赋值环谱}，$A_K$ 为其\Term{fraction-field}{分式域}上的\Term{abelian-scheme}{阿贝尔概形}，$\ell$ 为一个\Term{prime-number}{素数}。当 $\ell=p$ 时，假设 $K$ 的\Term{characteristic}{特征}为零（参见 5.11）。

 a) $A$ 在 $S$ 上具有\Term{good-reduction}{好约化}（2.2.9）的充要条件，是 $T_\ell(A_K)$ 在 $S$ 上具有\Term{good-reduction}{好约化}（5.12）。

 b) $A_K$ 在 $S$ 上具有\Term{potential-good-reduction}{本质好约化}，即存在 $K$ 的\Term{finite}{有限}\Term{extension}{扩张} $K'$，使 $A_{K'}$ 在 $S$ 于 $K'$ 中的\Term{normalization}{正规化} $S'$ 上具有\Term{good-reduction}{好约化}，其充要条件是 $T_\ell(A_K)$ 在 $S$ 上具有\Term{potential-good-reduction}{本质好约化}（5.12）。此时 $K'$ 可取为 $K$ 的\Term{finite}{有限}\Term{separable-extension}{可分扩张}。

 c) 以下条件\Term{equivalence}{等价}：
\begin{enumerate}
\item $A$ 在 $S$ 上具有\Term{semistable-reduction}{半稳定约化}（3.4）。
\item $T_\ell(A_K)$ 在 $S$ 上具有\Term{virtual-good-reduction}{阶数为 $2$ 的虚好约化}（5.12）。
\item $T_\ell(A_K)$ 在 $S$ 上具有\Term{virtual-good-reduction}{虚好约化}。
\end{enumerate}
\end{SGAresult}
\begin{proof}
a) 正是 5.10，列出只是为了备忘。

b) 第一个断言由 a) 和定义立即得到。要证明 b) 的最后一个断言，可假设 $\ell\ne p$。但由 a)，此时只需取一个\Term{finite}{有限}\Term{separable-extension}{可分扩张}
\SourcePage{398}
$K'$，对应于 $\pi=\operatorname{Gal}(\overline K/K)$ 的一个\Term{open-subgroup}{开子群} $\pi'$，使 $\pi'$ 与\Term{inertia-group}{惯性群} $I$ 的交包含在 $I$ 在 $T_\ell(A_K(\overline K))$ 上的\Term{representation}{表示}的\Term{kernel}{核}中；根据假设，这个\Term{kernel}{核}是开的（参见 5.12.1）。

c) 已经指出 $1)\Rightarrow2)$（5.12.2），而 $2)\Rightarrow3)$ 显然，所以只须证明 $3)\Rightarrow1)$。当 $\ell\ne p$ 时，这正是\Term{semistable-reduction}{半稳定约化}的 Galois 判据 3.5。当 $\ell=p$ 时，可如下进行。选一个\Term{prime-number}{素数} $\ell'\ne p$，对它应用 4.1 的讨论。设 $S$ \Term{strictly-local}{严格局部}，这是允许的（3.4.0）；取 $K$ 的\Term{finite}{有限} Galois \Term{extension}{扩张} $K'$，其群为 $G$，使 $A_{K'}$ 具有\Term{semistable-reduction}{半稳定约化}。需要证明 $G$ 在 (4.1.1) 出现的三个商上的\Term{representation}{表示}平凡；\Term{equivalence}{等价}地，它们的和\Term{representation}{表示}平凡；又\Term{equivalence}{等价}于后者的迹“平凡”，即为常数。4.3 a) 已经说明，此迹函数与 $\ell'\ne p$ 的选取无关；而下文 6.4 将给出一种也能用 $T_p(A_K)$ 来计算此函数的方法。当 $T_p(A_K)$ 具有\Term{virtual-good-reduction}{虚好约化}时，这个计算显然给出\Term{constant}{常值}函数。证明完成。
\end{proof}

\textbf{注记 5.13.1。}下面的说法似乎成立：若 $X_\eta$ 为 $\eta$ 上的 \Term{barsotti-tate-group}{Barsotti--Tate pro-$\ell$ 群}，在 $S$ 上具有\Term{virtual-good-reduction}{虚好约化}，并具有\Term{potential-virtual-good-reduction}{阶数为 $n$ 的本质虚好约化}，则它具有\Term{virtual-good-reduction}{阶数为 $n$ 的虚好约化}。（$\ell\ne p$ 时确实成立。）

\SourcePage{399}
下述结果是 5.7 和\Term{semistable-reduction}{半稳定约化}\Term{theorem}{定理} 3.6 的一个容易的推论，由 J.-P. Serre 向我指出。
\begin{SGAresult}{命题}{5.14（\Term{potential-good-reduction}{本质好约化}判据）}
设 $S$ 为\Term{function-field}{函数域}是 $K$ 的\Term{trait}{离散赋值环谱}，$A_K$ 为 $K$ 上的\Term{abelian-scheme}{阿贝尔概形}，$\ell$ 为不同于 $K$ 的\Term{characteristic}{特征}的\Term{prime-number}{素数}。

 a) 设 $\ell\ne p=\operatorname{char}k$。则下述条件\Term{equivalence}{等价}：
\begin{enumerate}
\item[(i)] $A_K$ 在 $S$ 上具有\Term{potential-good-reduction}{本质好约化}，即存在 $K$ 的\Term{finite}{有限}\Term{separable-extension}{可分扩张} $K'$，使 $A_{K'}$ 相对于 $S$ 在 $K'$ 中的\Term{normalization}{正规化} $S'$ 具有\Term{good-reduction}{好约化}。
\item[(ii)] \Term{inertia-group}{惯性群} $I$ 在 $U_\ell=T_\ell(A(\overline K))$ 上的\Term{action}{作用}本质平凡，即在 $I$ 的一个\Term{open-subgroup}{开子群}上平凡；也就是说，此\Term{action}{作用}通过 $I$ 的某个\Term{finite}{有限}\Term{quotient-group}{商群} $\Gamma$ 进行。
\item[(iii)] \Term{inertia-group}{惯性群} $I$ 在 $U_\ell\otimes_{\mathbb Z_\ell}\mathbb Q_\ell$ 上的\Term{action}{作用}是\Term{semisimple}{半单}的。
\end{enumerate}
 b) 设 $\ell=p$；若 $n=\dim A_K\ne0$，则 $I$ 在 $U_\ell$ 的\Term{automorphism-group}{自同构群}中的像不是\Term{finite-group}{有限群}。更确切地说，它在 $\bigwedge^{2n}U_\ell$ 的\Term{automorphism-group}{自同构群} $\mathbb Z_\ell^*$ 中的像具有\Term{finite-index}{有限指数}。$I$ 在 $U_\ell\otimes_{\mathbb Z_\ell}\mathbb Q_\ell$ 上的\Term{action}{作用}\Term{semisimple}{半单}，是 $A_K$ 具有\Term{potential-good-reduction}{本质好约化}的充分条件（但绝非必要条件）。
\end{SGAresult}
\begin{proof}
当 $\ell\ne p$ 时，(i) 与 (ii) 的\Term{equivalence}{等价}是\Term{good-reduction}{好约化}判据 2.2.9 的直接推论；由于 $\mathbb Q_\ell$ 的\Term{characteristic}{特征}为零，(ii) 蕴含 (iii)。\Term{homomorphism}{同态}
\[
\pi=\operatorname{Gal}(\overline K/K)\longrightarrow\mathbb Z_\ell^*
\simeq\operatorname{Aut}\bigl(\textstyle\bigwedge^{2n}U_\ell\bigr)
\]
等于 $\chi^n$，其中 $\chi$ 由
\SourcePage{400}
$ I$ 在 $\ell^\nu$ 次\Term{root-of-unity}{单位根}上的\Term{action}{作用}所给的\Term{character}{特征标}
\[
\mathbb Z_\ell^*=\operatorname{Aut}(\mathbb Z_\ell(1))
\]
诱导。已知，当 $\ell=p$ 时，$I$ 经此\Term{character}{特征标}的像确有\Term{finite-index}{有限指数} [33]。即使 $A_K$ 为具有\Term{good-reduction}{好约化}的\Term{elliptic-curve}{椭圆曲线}，$I$ 在 $U_p\otimes\mathbb Q_p$ 上的\Term{action}{作用}仍可能并非\Term{semisimple}{半单}，见 [28，IV A 2.4]。现在只须证明，对任意 $\ell$，(iii) 蕴含 (i)。为此，应用\Term{semistable-reduction}{半稳定约化}\Term{theorem}{定理} 3.6，便立即归结到 $A_K$ 在 $S$ 上具有\Term{semistable-reduction}{半稳定约化}的情形，因为把 $K$ 换成\Term{finite}{有限}\Term{separable-extension}{可分扩张} $K'$，即把 $I$ 换成\Term{open-subgroup}{开子群} $I'$，假设 (iii) 仍然保持（\SGAsourcegap{原印在此\Term{lie-bracket}{括号}中留白，缺说明或引注}）。当 $\ell\ne p$ 时，$I$ 在 $U_\ell\otimes\mathbb Q_\ell$ 上的\Term{action}{作用}是\Term{unipotent}{幺幂}的（3.5 (v)），又是\Term{semisimple}{半单}的，因而平凡，再用\Term{good-reduction}{好约化}判据 2.2.9 即得结论。下面给出一个对 $\ell$ 无须\Term{restriction}{限制}的证明，并同时得到略为精确的结果。
\end{proof}

\begin{SGAresult}{推论}{5.15}
沿用 5.14 的记号，设 $A_K$ 在 $S$ 上具有\Term{semistable-reduction}{半稳定约化}，且\Term{fixed-part}{固定部分}
\[
(U_{\ell,\mathbb Q_\ell})^f\overset{\mathrm{def}}=(U_\ell^f)_{\mathbb Q_\ell}
\]
在 $U_{\ell,\mathbb Q_\ell}$ 中有一个对 $I$ \Term{stable}{稳定}的\Term{complement}{补空间}。则 $A_K$ 具有\Term{good-reduction}{好约化}。
\end{SGAresult}
\begin{proof}
只需证明在这些条件下有 $U_\ell=U_\ell^f$；由 5.6，这意味着 $A_o^o$ 的\Term{toric-part}{环面部分} $T_o$ 为零，故 $A_o^o$ 为\Term{abelian-scheme}{阿贝尔概形}，从而 $A$ 确为\Term{abelian-scheme}{阿贝尔概形}（2.2.9）。显然可以假设 $k$ \Term{separably-closed}{可分闭}。于是，5.15 所要求的\Term{complement}{补空间}与 $U_\ell$ 的交
\SourcePage{401}
定义了 $T_\ell(A_K)$ 的一个子 pro-$\ell$ 群 $F_K$，它与余\Term{toric-part}{环面部分} $T_\ell(A_K)/T_\ell(A_K)^f$ \Term{isogenous}{同源}。由 5.6，后者是\Term{unramified}{非分歧} pro-\Term{etale}{平展}群，因此 $F$ 也是如此。由 5.7，得到 $F_K$ 包含于 $T_\ell(A_K)^f$，所以 $F_K=0$，这就证明 $U_\ell$ 等于其\Term{fixed-part}{固定部分}。
\end{proof}

\SGAsection{6}{非半稳定约化情形中关于构造 $T_p(A)$ 的环面部分与固定部分的注记}
\textbf{6.1。}5.1 中给出的 $T_\ell(A_K)$ 的\Term{toric-part}{环面部分}的构造，只要 $S$ \Term{complete}{完备}就有意义，无须假设 $A^o$ 为 $\ell$ \Term{divisible}{可除}的；即在 $p=\ell$ 的情形，无须假设具有\Term{semistable-reduction}{半稳定约化}（这个假设此前只是为表述\Term{orthogonality-theorem}{正交性定理} 5.2 而使用）。其实，只假设 $S$ 是 \Term{henselian}{Hensel} 的，也能构造 $T_p(A_K)^t$：证明在\Term{completion}{完备化}上的\Term{toric-part}{环面部分} $T_p(A_{\widehat K})^t$，即用 5.1 的方法构造的部分，来自 $T_p(A_K)$ 的一个 \Term{barsotti-tate-group}{Barsotti--Tate 子 pro-群}即可（显然它唯一确定）。更确切地说，可以在 $S$ 上构造一个 \Term{barsotti-tate-group}{Barsotti--Tate $p$-pro-群} $T_p(A)^t$，及一个\Term{projective-system}{投射系}的\Term{homomorphism}{同态} $T_p(A)^t\to T_p(A)$，它逐项为\Term{monomorphism}{单态射}，而在 $\widehat S$ 上取\Term{inverse-image}{逆像}时，给出用 5.1 的方法构造的相应\Term{object}{对象}。仿照 5.1 的构造即可看出这一点，但不再\Term{lifting}{提升} $A_o$ 的整个极大\Term{torus}{环面} $T_o$，而对每个整数 $\nu\ge0$，只把\Term{subgroup}{子群} ${}_{p^\nu}T_o$ \Term{lifting}{提升}为 ${}_{p^\nu}A^o$ 的一个\Term{multiplicative-type}{乘法型}\Term{subgroup}{子群}
\SourcePage{402}
$({}_{p^\nu}A^o)^t$。关于这种\Term{lifting}{提升}的存在性与唯一性，注意\Term{functor}{函子}
\[
(\mathrm{Sch})^o_{/S}\longrightarrow\mathrm{Ens},\qquad
S'\longmapsto\{A^o_{S'}\text{ 的有限乘法型子群}\}
\]
可由一个 $S$ 上\Term{finite-type}{有限型}\Term{scheme}{概形} $X$ \Term{representable}{表示}（参见 SGA 3 XI 3.12 a) 的证明）；由 SGA 3 IX 3.6 bis，$X$ 在 $S$ 上\Term{etale}{平展}。因此，$X_s$ 在 $s$ 上的每个给定\Term{section}{截面}，恰有一个 $X$ 在 $S$ 上的\Term{section}{截面}延拓它。这就证明了我们的断言。

更\Term{generic}{一般}地，这一论证说明：对 $A_o$ 的每个\Term{finite}{有限}\Term{multiplicative-type}{乘法型}\Term{subgroup}{子群} $H_o$，存在唯一的 $A$ 的\Term{finite}{有限}\Term{multiplicative-type}{乘法型}\Term{subgroup}{子群} $H$ 延拓它。因此，对 $A_o$ 的每个\Term{toroidal}{环面型}（toroïdal，即各项均为\Term{multiplicative-type}{乘法型}）\Term{barsotti-tate-group}{Barsotti--Tate 子 $p$-群} $H_o$，存在唯一的 $A$ 的 \Term{barsotti-tate-group}{Barsotti--Tate 子 $p$-群} $H$ 延拓它；因而它给出 ${}_{p^\infty}A_K$ 的一个 \Term{barsotti-tate-group}{Barsotti--Tate 子 $p$-群}，或者说，给出 $T_p(A_K)$ 的一个 \Term{barsotti-tate-group}{Barsotti--Tate 子 $p$-pro-群}。此过程特别适用于 $H_o$ 为 $A_o$ 中最大的\Term{toroidal}{环面型} \Term{barsotti-tate-group}{Barsotti--Tate 子 $p$-群}的情形：先取 $A_o$ 的最大 $p$ \Term{divisible}{可除}\Term{subgroup}{子群} $\overline A_o$（对充分大的 $N$，它等于 $p^N A_o^o$），再取 ${}_{p^\infty}\overline A_o$ 的“\Term{toroidal-part}{环面型部分}”（例如，可用 SGA 3 XVII 7.2.1 定义）。这样得到 ${}_{p^\infty}A_K$ 的一个 \Term{barsotti-tate-group}{Barsotti--Tate 子群}，可以称为后者的\Term{toroidal-part}{环面型部分}。

\textbf{6.2。}重新取 $A_o$ 的最大 $p$ \Term{divisible}{可除}\Term{subgroup-scheme}{子群概形} $\overline A_o$，并考虑 \Term{barsotti-tate-group}{Barsotti--Tate $p$-群} $F_o={}_{p^\infty}(\overline A_o)$；它显然可刻画为 $A_o$（或 $A_o^o$）中最大的 \Term{barsotti-tate-group}{Barsotti--Tate 子 $p$-群}。自然会问，能否将此\Term{barsotti-tate-group}{Barsotti--Tate 群}
\SourcePage{403}
延拓为 $A$ 的一个 \Term{barsotti-tate-group}{Barsotti--Tate 子群} $F$，使它起到 $A^o$ 为 $\ell$ \Term{divisible}{可除}时 (2.2.3.2) 的\Term{fixed-part}{固定部分} $T_\ell(A)^f$ 所起的\Term{action}{作用}。在\Term{semistable-reduction}{半稳定}情形，$\overline A_o=A_o^o$，由 2.2 的方法可肯定回答这个问题。在\Term{generic}{一般}情形，若所考虑的问题有解，则容易看出解唯一；而 $T_p(A_K)$ 与 $T_p(A'_K)$ 中对应的子 pro-群（这时仍值得称为相应 pro-群的\Term{fixed-part}{固定部分}）仍满足\Term{orthogonality-theorem}{正交性定理} 5.2；最后，\Term{fixed-part}{固定部分} $T_p(A_K)^f$ 仍可如 5.8 所述那样刻画。

\textbf{6.2.1。}不过，我们将说明，\Term{generic}{一般}而言上述问题没有肯定答案，也就是说，\Term{generic}{一般}无法定义\Term{fixed-part}{固定部分} $T_p(A)^f$（在非\Term{semistable-reduction}{半稳定约化}情形）。事实上，根据 Serre--Tate 的一个\Term{lifting}{提升}\Term{theorem}{定理} [30]，当 $S$ \Term{complete}{完备}时，\Term{fixed-part}{固定部分}的存在性，\Term{equivalence}{等价}于在各 $A_n$ 内存在 $\overline A_o$ 的\Term{infinitesimal-lifting}{无穷小提升} $\overline A_n$；又\Term{equivalence}{等价}于存在 $S$ 上的\Term{smooth}{光滑}\Term{group-scheme}{群概形} $\overline A$，它是一个\Term{abelian-scheme}{阿贝尔概形}被\Term{torus}{环面} $T$ 的\Term{group-extension}{扩张}，其\Term{special-fiber}{特殊纤维}为 $\overline A_o$，并存在一个\Term{formal-homomorphism}{形式同态} $\widehat{\overline A}\to\widehat A$，延拓\Term{special-fiber}{特殊纤维}上的包含。因此，当 $\overline A_o$ 为\Term{abelian-scheme}{阿贝尔概形}，即 $T_o=0$ 时，这就\Term{equivalence}{等价}于存在一个延拓 $\overline A_o$ 的\Term{abelian-scheme}{阿贝尔概形} $\overline A$，以及一个延拓\Term{special-fiber}{特殊纤维}的包含的\Term{homomorphism}{同态} $\overline A\to A$。因此，只须给出一个情形，使 $T_o=0$，$A_o^o$ 既非\Term{unipotent}{幺幂}又非\Term{proper}{固有}，且 $A_K$ 在 $K$ 上\Term{simple}{单}。

\SourcePage{404}
\textbf{6.2.2。}下面是 M. Raynaud 的一个例子，依据如下观察：若 $S$ 为\Term{trait}{离散赋值环谱}，$S'$ 为 $S$ 在 $K$ 的一个\Term{finite}{有限}\Term{separable-extension}{可分扩张} $K'$ 中的\Term{normalization}{正规化}，$A'$ 为 $S'$ 上的\Term{abelian-scheme}{阿贝尔概形}，则 $S$ 上的\Term{group-scheme}{群概形}
\[
A=\prod_{S'/S}A'/S'
\]
具有 Néron 性质，即它等同于其\Term{generic-fiber}{一般纤维}的 \Term{neron-model}{Néron 模型}（该\Term{generic-fiber}{一般纤维}确为\Term{abelian-scheme}{阿贝尔概形}）。由定义，以及 $A'$ 本身在 $S'$ 上具有 Néron 性质这一事实，这一点立即可得。于是
\[
A_o=\prod_{S'_o/S_o}(A'\times_{S'}S'_o)/S'_o
\]
是一个\Term{abelian-scheme}{阿贝尔概形}被\Term{connected}{连通}\Term{unipotent}{幺幂}\Term{scheme}{概形}的\Term{group-extension}{扩张}；若 $A'\ne0$，前者非零；若 $A'\ne0$ 且 $S'$ 在 $S$ 上\Term{ramified}{分歧}，后者也非零。因此，只须给出一个例子，使 $A'\ne0$，$S'$ 在 $S$ 上\Term{ramified}{分歧}（$S$ 为\Term{complete}{完备}\Term{local-scheme}{局部}\Term{trait}{离散赋值环谱}），并使
\[
A_K=\prod_{K'/K}A'_{K'}/K'
\]
在 $K'$ 上\Term{simple}{单}。取 $K'$ 为 $S$ 的一个\Term{ramified}{分歧}二次\Term{extension}{扩张}（这样的\Term{extension}{扩张}总是存在），取 $A'$ 为 $S'$ 上的一条\Term{elliptic-curve}{椭圆曲线}。于是 $A$ 的\Term{relative-dimension}{相对维数}为 $2$，而 $A_K$ 在 $K$ 上\Term{simple}{单}，是说不存在非零\Term{homomorphism}{同态} $B_K\to A_K$，其中 $B_K$ 为 $K$ 上\Term{dimension}{维数}是 $1$ 的\Term{abelian-scheme}{阿贝尔概形}。给出这样的\Term{homomorphism}{同态}，\Term{equivalence}{等价}于给出非零\Term{homomorphism}{同态} $B_{K'}\to A'_{K'}$；后者必为\Term{isogeny}{同源}。因此，$A_K$ 在 $K$ 上\Term{simple}{单}，\Term{equivalence}{等价}于 $A'_{K'}$ 不与任何来自 $K$ 的\Term{elliptic-curve}{椭圆曲线}\Term{isogenous}{同源}。若反之，则 $A'$ 的\Term{j-invariant}{$j$ 不变量} $j$ 与它在 $K'$ 的非平凡 $K$-\Term{automorphism}{自同构}下的\Term{conjugacy}{共轭} $\overline j$，在\Term{prime-field}{素域}上\Term{algebraically-dependent}{代数相关}；若选取 $j$，使 $a=j+\overline j=\operatorname{Tr}_{K'/K}j$ 在\Term{prime-field}{素域}上\Term{transcendental}{超越}，那么这也\Term{equivalence}{等价}于 $j$ 与 $a$ \Term{algebraically-dependent}{代数相关}。由\Term{cardinal}{基数}的原因，存在元素
\SourcePage{405}
$a\in V$，在 $K$ 的\Term{prime-field}{素子域} $P$ 上\Term{transcendental}{超越}；进而存在 $V'$（$V$ 在 $K'$ 中的\Term{normalization}{正规化}）中的元素 $b$，在 $P(a)$ 上\Term{transcendental}{超越}，且 $\operatorname{Tr}_{V'/V}b=0$。只须取 $j=a+b\in V'$。此外，可以要求 $b\in\mathfrak r(V')$，且 $a$ 在 $k=V/\mathfrak m$ 中的像为 $k$ 的一个给定元素，选它不同于\Term{exceptional-value}{例外值} $j_o=0$ 和 $j_o=12^3$（注意，即使 $k$ 只有两个元素也能如此，因为这时上述两个 $j_o$ 值相同）。于是 $j$ 在 $k(s')$ 中的像 $j_o$ 不等于 $0,12^3$，已知这时确实存在 $S'$ 上以 $j$ 为\Term{j-invariant}{$j$ 不变量}的\Term{elliptic-curve}{椭圆曲线}。

\textbf{6.3。}现在回到仅假设 $S$ 为 \Term{henselian}{Hensel} 的情形。类似于 $\ell\ne p$ 时 4.1 的定义，此时可以给出 $T_p(A_K)$ 的一个\Term{filtration}{滤过}
\[
T_p(A_K)\supset T_p(A_K)^{\mathrm{ef}}\supset T_p(A_K)^{\mathrm{et}}\supset0,\tag{6.3.1}
\]
它由（\Term{strict-subgroup}{严格}的）\Term{barsotti-tate-group}{Barsotti--Tate 子 pro-群}组成；这两个部分分别称为 $T_p(A_K)$ 的\Term{essentially-fixed-part}{本质固定部分}和\Term{essentially-toric-part}{本质环面部分}。为此，选一个 $K$ 的\Term{finite}{有限} Galois \Term{extension}{扩张} $K'$，使 $A_{K'}$ 在 $S$ 于 $K'$ 中的\Term{normalization}{正规化} $S'$ 上具有\Term{semistable-reduction}{半稳定约化}。于是 $T_p(A_{K'})=(T_p(A_K))_{K'}$ 的\Term{fixed-part}{固定部分} $T_p(A_{K'})^f$ 有定义（2.2.3），\Term{toric-part}{环面部分}也有定义（6.1）；后者又恰好是 $T_p(A'_{K'})$ 的\Term{fixed-part}{固定部分} $T_p(A'_{K'})^f$ 的\Term{orthogonal-complement}{正交补}（5.2）。显然有 $T_p(A_{K'})^t\subset T_p(A_{K'})^f$（由 5.3 可得，或直接由 6.1 中\Term{toric-part}{环面部分}的构造得到），且 $K'/K$ 的 Galois 群保持 $(T_p(A_K))_{K'}$ 的这两个\Term{strict-subgroup}{严格}子 pro-群\Term{invariant}{不变}。

\SourcePage{406}
因此，通过\Term{descent}{下降}便得到所要的\Term{filtration}{滤过} (6.3.1)。它与 $K'$ 的选取无关，这由\Term{semistable-reduction}{半稳定约化}的性质 3.3 得到。此外，通过\Term{descent}{下降}显然有\Term{orthogonality}{正交性}关系
\[
T_p(A_K)^{\mathrm{et}}=(T_p(A'_K)^{\mathrm{ef}})^\perp,\qquad
T_p(A_K)^{\mathrm{ef}}=(T_p(A'_K)^{\mathrm{et}})^\perp.\tag{6.3.2}
\]
可以本着 5.8 的精神，仅用 $S$ 的\Term{fraction-field}{分式域} $K$ 上的 \Term{barsotti-tate-group}{Barsotti--Tate pro-群} $X=T_p(A_K)$，内在地刻画\Term{filtration}{滤过} (6.3.1)，至少目前在\Term{mixed-characteristic}{混合特征}情形可如此：\Term{essentially-fixed-part}{本质固定部分}是“在 $S$ 上具有\Term{potential-good-reduction}{本质好约化}”的最大 \Term{barsotti-tate-group}{Barsotti--Tate 子 pro-群}；也就是说，在 $K$ 的一个\Term{finite}{有限}\Term{separable-extension}{可分扩张} $K'$ 之后，所得 \Term{barsotti-tate-group}{Barsotti--Tate pro-群} $X_{K'}$ 在 $S$ 于 $K'$ 中的\Term{normalization}{正规化} $S'$ 上具有\Term{good-reduction}{好约化}，即延拓为 $S'$ 上的 \Term{barsotti-tate-group}{Barsotti--Tate pro-群}（由已引用的 Tate \Term{theorem}{定理} [33，\Term{theorem}{定理} 4]，在唯一\Term{isomorphism}{同构}的意义下必然唯一）。然后，通过\Term{orthogonality}{正交性} (6.3.2)，由\Term{essentially-fixed-part}{本质固定部分}定义\Term{essentially-toric-part}{本质环面部分}。

再指出 5.8 与定义的下述直接推论；当 $\ell=p$ 时，它其实归结为 5.12 a)。
\begin{SGAresult}{命题}{6.3.3}
设 $\ell$ 为一个\Term{prime-number}{素数}。当 $\ell=p$ 时，假设 $\operatorname{char}K=0$\footnote{这个假设大概并无必要，参见 5.11。}。$A_K$ 在 $S$ 上具有\Term{potential-good-reduction}{本质好约化}（5.12 a)）的充要条件，是\Term{essentially-fixed-part}{本质固定部分} $T_\ell(A_K)^{\mathrm{ef}}$（4.1.1 及 6.3.1）等于整个 $T_\ell(A_K)$。
\end{SGAresult}

\SourcePage{407}
\textbf{6.4。}注意，当 $\ell\ne p$ 时，我们在 4.3 a) 给出了 $S$ 的\Term{inertia-group}{惯性群} $I$ 上一个引人注意的整值函数；它与 $\ell$ 的选取无关。当\Term{residue-field}{剩余域}为\Term{perfect-field}{完美域}时，只要知道此函数，就足以恢复 $A$ 的\Term{conductor-exponent}{导子指数}（4.5）：在 $I$ 的一个充分大的商 $\Gamma=I/I'$ 上，取上述函数与 $\Gamma$ 上的 Artin \Term{conductor}{导子}的内积即可（参见 (4.6.2)）。值得注意的是，此 $I$ 上的整值函数也能由 $\ell=p$ 时的 $T_\ell(A_K)$ 恢复——至少目前在\Term{mixed-characteristic}{混合特征}时可以。当\Term{residue-field}{剩余域} $k$ 为\Term{finite-field}{有限域}时，类似的观察也适用于 4.3 b) 所考虑的取值于 $\mathbb Q$ 的函数：它定义在 $\pi$ 的\Term{subgroup}{子群} $\pi^!$ 上，该\Term{subgroup}{子群}由那些在 $\pi_o$ 中的像形如 $f_q^n$（$n\in\mathbb Z$）的元素组成。

下面说明这些断言的证明原理。使用\Term{filtration}{滤过} (6.3.1)，问题归结为：给定 $K$ 上一个“在 $S$ 上具有\Term{potential-good-reduction}{本质好约化}”的 \Term{barsotti-tate-group}{Barsotti--Tate $p$-pro-群} $Y$（含义如 6.5），为它赋予一个适当的函数，定义在 $I$ 上，相应地定义在 $\pi^!$ 上。（$Y$ 是\Term{filtration}{滤过} (6.3.1) 给出的三个商之一。）在 4.3 a) 的情形，作 $S$ 的适当\Term{residue-extension}{剩余扩张}（EGA $0_{\mathrm{III}}$ 10.3.1），便归结到\Term{residue-field}{剩余域}为\Term{perfect-field}{完美域}（若愿意，甚至可为\Term{algebraically-closed}{代数闭域}）的情形。取 $K$ 的\Term{finite}{有限} Galois \Term{extension}{扩张} $K'$，使 $Y_{K'}$ 在 $S'$ 上具有\Term{good-reduction}{好约化}；于是有 $S'$ 上的 \Term{barsotti-tate-group}{Barsotti--Tate 群} $Y_{S'}$，以及 $G$ 在 $Y_{S'}$ 上与其在 $S'$ 上的\Term{action}{作用}相容的\Term{action}{作用}。若 $k'$ 为 $S'$ 的\Term{residue-field}{剩余域}，则 $G$
\SourcePage{408}
在 $k'$ 上\Term{action}{作用}（未必\Term{faithful-action}{忠实}，其\Term{invariant}{不变}子域为 $k$），并在 $Y_{k'}$ 上以与其在 $k'$ 上的\Term{action}{作用}相容的方式\Term{action}{作用}。由 Dieudonné 理论 [16 bis]，\Term{barsotti-tate-group}{Barsotti--Tate $p$-pro-群} $Y_{k'}$ \Term{canonical}{典范}地对应于 $k'$ 上的\Term{witt-vectors}{Witt 向量}环 $W(k')$ 上的一个\Term{finite-free-module}{有限型自由模} $M$（其\Term{rank}{秩}等于 $Y_{k'}$ 的\Term{rank}{秩}）；$G$ 显然在 $M$ 上\Term{action}{作用}，并与它在 $W(k')$ 上的\Term{action}{作用}相容（后者来自它在 $k'$ 上的\Term{action}{作用}）。由此得到 $G$ 的\Term{inertia-subgroup}{惯性子群} $\Gamma$ 在 $M$ 上的 $W(k')$-\Term{linear-action}{线性作用}；$\Gamma$ 是\Term{inertia-group}{惯性群} $I$ 的一个商。因此在 $\Gamma$ 上，进而在 $I$ 上，得到一个迹函数，这正是所要的函数。先验地它取值于 $W(k')$\footnote{事实上，由于此函数在 $\operatorname{Gal}(k'/k)$ 下\Term{invariant}{不变}，立即可知它总是取值于 $W(k)$。}。不过，当 $G$ 在 $Y_{S'}$ 上的\Term{action}{作用}来自它在 $S'$ 上的一个\Term{smooth}{光滑}且 $p$ \Term{divisible}{可除}的\Term{group-scheme}{群概形} $C'$ 上的\Term{action}{作用}，且 $Y=T_p(C')$ 时，此函数与用 $\ell\ne p$ 的 $T_\ell$ 定义的类似函数一致（参见 [25 bis，\Term{theorem}{定理} 7，第 737 页]），故在这种情形它取整数值。在 4.3 b) 的情形也类似：采用 4.3.2 证明中的方法，将位于某个 $f_q^n$（$n\ge0$）之上的元素 $g\in\pi$ 在 $M$ 上定义的\Term{action}{作用}“\Term{linear}{线性}化”，就能在 $\pi^!$ 上定义迹函数。
